\scanpage{281}
\mspage{22}

où le deuxième $\mathrm{Hom}$ est considéré comme un $\mathrm{Hom}$ d'objets dans $\mathrm{Cat}^{\unsure{\times|\vee}}=\underline{\mathrm{Hom}}((\mathrm{Cat}),\mathrm{Ens})$.

Ainsi, dans les deux exemples plus haut, où $\tau=\mathbb{\Delta}_2^+$ et $\tau'=\mathbb{\Delta}_1^+$ ou $\mathbb{\Delta}_2^+$, les foncteurs en sens inverse sont
\diagram{\begin{tikzcd}[ampersand replacement=\&,row sep=small,column sep=small]
\mathbb{\Delta}_2^+\quad\{ \& 0 \arrow[r] \& 1 \arrow[r] \& 2 \& \unsure{/}\\
\mathbb{\Delta}_1^+\quad\{ \& 0 \arrow[u,maps to] \arrow[rr] \&{} \& 1 \arrow[u,maps to] \&{}
\end{tikzcd}}
et
\diagram{\begin{tikzcd}[ampersand replacement=\&,row sep=small,column sep=small]
\mathbb{\Delta}_2^+\quad\{ \&0 \arrow[r] \& 1 \arrow[r,dash] \& 2\\
\mathbb{\Delta}_2^+\quad\{ \&0 \arrow[r] \arrow[u,maps to] \& 1 \arrow[r] \arrow[ur,maps to] \& 2 \arrow[u,maps to]
\end{tikzcd}}.

C'est là un principe général pour ramener toutes les ``opérations \inserted{$\pm$} canoniques sur des chemins \inserted{dans des cat. arbitraires}, et des énoncés les concernant, à des questions où $X$ disparaît, et où on travaille dans les catégories intervalles elles-mêmes.

Un des avantages techniques, c'est que ce sont des catégories fort petites, donnant lieu de plus à une intuition topologique \inserted{sûre} (qui risque de se volatiliser \inserted{par moments} en travaillant \deleted{\unsure{dans}} dans des catégories de chemins), et aussi que dans \unsure{de|des} telles catégories, \deleted{et les} \inserted{donc aussi \unsure{des|les}} catégories de foncteurs d'entre elles, il y a tjrs au plus
% archive sheet number: 280
